
%%FILE blueprint/src/chapter/char_fun.tex
\chapter{The characteristic function}

\section{Definition and first properties}

\begin{definition}[Characteristic function]\label{def:charFun}
\lean{MeasureTheory.charFun} \mathlibok
Let $\mu$ be a measure on a real inner product space $E$. The characteristic function of $\mu$, denoted by $\hat{\mu}$, is the function $E \to \mathbb{C}$ defined by
\begin{align*}
\hat{\mu}(t) = \int_x e^{i \langle t, x \rangle} d\mu(x) \: .
\end{align*}
The characteristic function of a random variable $X$ is defined as the characteristic function of $\mathcal L(X)$.
\end{definition}


\begin{lemma}\label{lem:charFun_bounded}
\lean{MeasureTheory.norm_charFun_le_one} \mathlibok
\uses{def:charFun}
For all $t$, $\Vert\hat{\mu}(t)\Vert \le 1$.
\end{lemma}

\begin{proof}\leanok
\end{proof}


\begin{lemma}\label{lem:charFun_contDiff}
\leanok
\lean{contDiff_charFun}
\uses{def:charFun}
If $\mu[|X|^n] < \infty$ then the characteristic function of $\mu$ is continuously differentiable of order $n$.
\end{lemma}

\begin{proof}\leanok
\end{proof}


\begin{lemma}\label{lem:charFun_continuous}
\leanok
\lean{continuous_charFun}
\uses{def:charFun}
The characteristic function is a continuous function.
\end{lemma}

\begin{proof}\leanok
\uses{lem:charFun_contDiff}
\end{proof}


\begin{lemma}\label{lem:charFun_neg}
\lean{MeasureTheory.charFun_neg} \mathlibok
\uses{def:charFun}
$\hat{\mu}(-t) = \overline{\hat{\mu}(t)}$.
\end{lemma}

\begin{proof}\leanok
\end{proof}


\begin{lemma}\label{lem:charFun_smul}
\mathlibok
\lean{MeasureTheory.charFun_map_smul}
\uses{def:charFun}
For $a \in \mathbb{R}$ and $X$ a random variable with law $\mu$, the characteristic function of $a X$ is $t \mapsto \hat{\mu}(at)$.
\end{lemma}

\begin{proof}\leanok
\end{proof}


\begin{lemma}\label{lem:charFun_add_of_indep}
\mathlibok
\lean{MeasureTheory.charFun_conv}
\uses{def:charFun}
If two random variables $X, Y : \Omega \to S$ with laws $\mu$ and $\nu$ are independent, then $X+Y$ has characteristic function $\hat{\mu} \times \hat{\nu}$.
\end{lemma}

\begin{proof}\leanok
\end{proof}


\begin{lemma}\label{lem:integral_exp_I}
\mathlibok
\lean{integral_exp_mul_I_eq_sin}
\begin{align*}
    \int_{-r}^r e^{i t} \, d t = 2 \sin r \: .
\end{align*}
\end{lemma}

\begin{proof}\leanok
\end{proof}


\begin{lemma}\label{lem:integral_charFun}
\uses{def:charFun}
\mathlibok
\lean{integral_exp_mul_I_eq_sinc}
\begin{align*}
    \int_{-r}^r \hat{\mu}(t) dt = 2 r \int \frac{\sin(r x)}{r x} \, d\mu(x)
\end{align*}
\end{lemma}

\begin{proof}
\uses{lem:integral_exp_I}
\leanok
We will use Fubini's theorem to exchange the order of integration.
\begin{align*}
    \int_{-r}^r \hat{\mu}(t) \, dt
    &= \int_{-r}^r \int e^{i t x} \, d\mu(x) \, dt
    \\
    &= \int \left( \int_{-r}^r e^{i t x} \, dt \right) \, d\mu(x)
    \\
    &= \int 2 \frac{\sin(r x)}{x} \, d\mu(x)
    \\
    &= 2 r \int \frac{\sin(r x)}{r x} \, d\mu(x)
\end{align*}
\end{proof}


% \begin{lemma}\label{lem:charFun_bound_small}
% \leanok
% \lean{ProbabilityTheory.measure_Icc_le_charFun}
% \uses{def:charFun}
% For $\mu$ a probability measure on $\mathbb{R}$ and $r > 0$,
% \begin{align*}
% \mu [-r, r]
% &\le 2 r \int_{-1/r}^{1/r} |\hat{\mu}(t)| dt
% \: .
% \end{align*}
% \end{lemma}

% \begin{proof}
% TODO
% \end{proof}


\begin{lemma}\label{lem:charFun_bound_large}
\mathlibok
\lean{MeasureTheory.measureReal_abs_gt_le_integral_charFun}
\uses{def:charFun}
For $\mu$ a probability measure on $\mathbb{R}$ and $r > 0$,
\begin{align*}
\mu \left\{x \mid |x| > r\right\}
&\le \frac{r}{2} \int_{-2/r}^{2/r} (1 - \hat{\mu}(t))dt
\: .
\end{align*}
\end{lemma}

\begin{proof}
\uses{lem:integral_charFun, lem:charFun_bounded}
\leanok
We will use Lemma~\ref{lem:integral_charFun} with $r$ replaced by $2/r$.
\begin{align*}
    \frac{r}{2} \int_{-2/r}^{2/r} (1 - \hat{\mu}(t))dt
    &= 2 - \frac{r}{2} \int_{-2/r}^{2/r} \hat{\mu}(t)dt
    \\
    &= 2 - 2 \int \frac{\sin(2 x / r)}{2 x / r} \, d\mu(x)
    \\
    &= 2 \int \left(1 - \frac{\sin(2 x / r)}{2 x / r}\right) \, d\mu(x)
    \\
    &\ge 2 \int_{|x| \ge r} \left(1 - \frac{\sin(2 x / r)}{2 x / r}\right) \, d\mu(x)
    \\
    &\ge \mu\{x \mid |x| > r\}
    \: .
\end{align*}
In the end we used that $\sin x \le x /2$ for $x \ge 2$.
\end{proof}


\begin{lemma}\label{lem:charFun_bound_inner}
\mathlibok
\lean{MeasureTheory.measureReal_abs_inner_gt_le_integral_charFun}
\uses{def:charFun}
For $\mu$ a probability measure on an inner product space $E$ and $r > 0$, $a \in E$,
\begin{align*}
\mu \left\{x \mid \vert\langle a, x\rangle\vert > r\right\}
&\le \frac{r}{2} \int_{-2/r}^{2/r} (1 - \hat{\mu}(t \cdot a))dt
\: .
\end{align*}
\end{lemma}

\begin{proof}\uses{lem:charFun_bound_large}
\leanok
TODO (see code)
\end{proof}

% \begin{lemma}\label{lem:abs_sub_charFun}
% \leanok
% \lean{ProbabilityTheory.abs_sub_charFun_le}
% \uses{def:charFun}
% Let $X$ be a random variable with law $\mu$. Then for any $s, t$,
% \begin{align*}
% \vert \hat{\mu}(s) - \hat{\mu}(t) \vert
% \le 2 \mu\left[ \left\vert (s - t) X\right\vert \wedge 1\right]
% \: .
% \end{align*}
% \end{lemma}

% \begin{proof}
% \begin{align*}
% \vert \hat{\mu}(s) - \hat{\mu}(t) \vert
% &= \vert \mu\left[e^{isX} - e^{itX}\right] \vert
% \\
% &\le \mu\left[\vert e^{isX} - e^{itX} \vert\right]
% \\
% &= \mu\left[\vert 1 - e^{i(t - s)X} \vert\right]
% \le 2 \mu\left[ \left\vert (s - t) X\right\vert \wedge 1\right]
% \: .
% \end{align*}
% \end{proof}


\section{Convergence of characteristic functions and weak convergence of measures}


\begin{lemma}\label{lem:tendsto_M_of_tendsto_charFun}
\leanok
\lean{MeasureTheory.ProbabilityMeasure.tendsto_charPoly_of_tendsto_charFun}
\uses{def:charFun}
Let $\mu, \mu_1, \mu_2, \ldots$ be probability measures on $\mathbb{R}^d$ with characteristic functions $\hat{\mu}, \hat{\mu}_1, \hat{\mu}_2, \ldots$ such that $\hat{\mu}_n \to \hat{\mu}$.
Then for all exponential polynomial $f \in \mathcal M$, $\mu_n[f] \to \nu[f]$.
\end{lemma}

\begin{proof}\uses{lem:starSubalgebra_expPoly}\leanok
\end{proof}


\begin{lemma}\label{lem:ext_charFun}
\mathlibok
\lean{MeasureTheory.Measure.ext_of_charFun}
Two finite measures on a complete separable metric space are equal iff they have the same characteristic function.
\end{lemma}

\begin{proof}\uses{thm:separating_starSubalgebra, lem:separates_points_expPoly}\leanok
See Mathlib PR \#19783.

The sub-algebra $\mathcal M$ of exponential polynomials is separating by Lemma~\ref{lem:separates_points_expPoly} and~\ref{thm:separating_starSubalgebra}.
Equality of characteristic functions implies equality on $\mathcal M$, which then implies equality of the measures since $\mathcal M$ is separating.
\end{proof}


% \begin{lemma}\label{lem:tight_iff_charFun_equicontinuous}
% \leanok
% \lean{isTightMeasureSet_iff_equicontinuousAt_charFun}
% \uses{def:tight}
% A family of measures $(\mu_a)$ on a finite dimensional inner product space is tight iff the family of functions $(\hat{\mu}_a)$ is equicontinuous at 0.
% \end{lemma}

% \begin{proof}\uses{lem:charFun_bound_small, lem:charFun_bound_inner, lem:abs_sub_charFun}
% TODO: equicontinuous implies tight.

% For the other direction, by Lemma~\ref{lem:abs_sub_charFun}, for all $a$,
% \begin{align*}
%     \vert \hat{\mu}_a(s) - \hat{\mu}_a(t) \vert
%     \le 2 \mu_a\left[ \left\vert (s - t) X\right\vert \wedge 1\right]
%     \: .
% \end{align*}
% Since $(\mu_a)$ is tight, the right-hand side converges to 0 as $s - t \to 0$ uniformly in $a$.
% We obtain equicontinuity at 0 of the family $(\hat{\mu}_a)$.
% \end{proof}


\begin{lemma}\label{lem:tight_of_tendsto_charFun}
\leanok
\lean{isTightMeasureSet_of_tendsto_charFun}
\uses{def:tight}
Let $(\mu_n)_{n \in \mathbb{N}}$ be measures on $\mathbb{R}^d$ with characteristic functions $(\hat{\mu}_n)$. If $\hat{\mu}_n$ converges pointwise to a function $f$ which is continuous at 0, then $(\mu_n)$ is tight.
\end{lemma}

\begin{proof}\uses{lem:charFun_bound_inner, lem:isTightMeasureSet_iff_basis, lem:charFun_bounded}
\leanok
By Lemma~\ref{lem:charFun_bound_inner} and dominated convergence, for all $a \in \mathbb{R}^d$ and $r > 0$,
\begin{align*}
    \lim_{n \to +\infty} \mu_n \left\{x \mid |\langle a, x\rangle| \ge r\right\}
    &\le \lim_{n \to +\infty} \frac{r}{2} \int_{-2/r}^{2/r} (1 - \hat{\mu}_n(\langle a, t\rangle))dt
    \\
    &= \frac{r}{2} \int_{-2/r}^{2/r} (1 - f(\langle a, t\rangle))dt
\end{align*}
The right-hand side converges to 0 as $r \to +\infty$ by continuity of $f$ at 0.
This shows that the family is tight.
\end{proof}

Remark: the finite dimension is necessary.
Counterexample: $\ell^2$ with $\mu_n$ the law of $X_n = \sum_{k=1}^n \zeta_k e_k$ where $e_k = (0, \ldots, 0, 1, 0 \ldots)$ and the $\zeta_k$ are i.i.d. $\mathcal N(0,1)$.
Then $\hat{\mu}_n(t) \to e^{- \Vert t \Vert^2 / 2}$ for all $t \in \ell^2$ but $(\mu_n)$ is not tight.


\begin{theorem}[Convergence of characteristic functions and weak convergence of measures]
\label{thm:charFun_tendsto_iff_measure_tendsto}
\leanok
\lean{MeasureTheory.ProbabilityMeasure.tendsto_iff_tendsto_charFun}
\uses{def:weak_cvg_measure}
Let $\mu, \mu_1, \mu_2, \ldots$ be probability measures on $\mathbb{R}^d$ with characteristic functions $\hat{\mu}, \hat{\mu}_1, \hat{\mu}_2, \ldots$. Then $\mu_n \xrightarrow{w} \mu$ iff for all $t$, $\hat{\mu}_n(t) \to \hat{\mu}(t)$.
\end{theorem}

\begin{proof}\uses{thm:separating_starSubalgebra, lem:separates_points_expPoly, lem:tight_of_tendsto_charFun, lem:cvg_of_separating, lem:tendsto_M_of_tendsto_charFun, lem:charFun_continuous}
\leanok
For all $t$, $x \mapsto e^{i \langle t, x \rangle}$ is a bounded continuous function. Hence by the definition of convergence in distribution, $\mu_n \xrightarrow{w} \mu \implies \hat{\mu}_n(t) \to \hat{\mu}(t)$ for all $t$.

Let's now prove the other direction.
If the characteristic functions converge pointwise, we get from Lemma~\ref{lem:tight_of_tendsto_charFun} that the family of measures is tight.
Then apply Lemma~\ref{lem:cvg_of_separating} to the subalgebra $\mathcal M$: it suffices to show that $\mu_n[f] \to \mu[f]$ for all $f \in \mathcal M$. This is obtained with Lemma~\ref{lem:tendsto_M_of_tendsto_charFun}.
\end{proof}

%%FILE blueprint/src/chapter/clt.tex
\chapter{The central limit theorem}

\section{The central limit theorem for real random variables}

\begin{lemma}\label{lem:deriv_charFun}
\leanok
\lean{iteratedDeriv_charFun}
\uses{def:charFun}
Let $X$ be a real random variable with characteristic function $\phi$ with $\mathbb{E}[\vert X \vert^n] < \infty$ and let $k \le n$.
Then the $k$\textsuperscript{th} derivative of $\phi$ is
$\phi^{(k)}(t) = \mathbb{E}[i X]^k e^{i t X}$ \: .
\end{lemma}

\begin{proof}\leanok
\uses{lem:charFun_contDiff}
\end{proof}

\begin{lemma}[Peano form of Taylor's theorem]\label{lem:taylor_peano}
\mathlibok
\lean{taylor_isLittleO}
For a $n$-times continuously differentiable function $f : \mathbb{R} \to E$, for any $x_0\in\mathbb{R}$,
\begin{align*}
f(x) - \sum_{i=0}^n\frac{f^{(i)}(x_0)}{i!}(x-x_0)^i = o((x - x_0)^k)
\: .
\end{align*}
\end{lemma}

\begin{proof}\leanok
\end{proof}

\begin{lemma}\label{lem:iIndepFun_iff_pi_map_eq_map}
\mathlibok
\lean{ProbabilityTheory.iIndepFun_iff_map_fun_eq_pi_map}
A finite collection of random variables are independent iff their joint law is the product of their respective laws.
\end{lemma}

\begin{proof}\leanok
\end{proof}

\begin{lemma}\label{lem:charFun_taylor}
\leanok
\lean{taylor_charFun}
Let $X$ be a real random variable with characteristic function $\phi$, with $\mathbb{E}[\vert X \vert^n] < \infty$. Then as $t \to 0$,
\begin{align*}
\phi(t) = \sum_{k=0}^n \frac{(it)^k \mathbb{E}[X^k]}{k!} + o(t^n)
\: .
\end{align*}
\end{lemma}

\begin{proof}\leanok
\uses{lem:deriv_charFun, lem:taylor_peano}
\end{proof}

\begin{lemma}\label{lem:tendsto_pow_exp_of_isLittleO}
\leanok
\lean{tendsto_pow_exp_of_isLittleO}
For $t\in\mathbb{C}$, $\lim_{n\to\infty}(1+t/n+o(1/n))^n=\exp(t)$ (where the little-$o$ term may be complex).
\end{lemma}

\begin{proof}\leanok
\end{proof}

\begin{theorem}[Central limit theorem]\label{thm:clt}
    \leanok
    \lean{ProbabilityTheory.central_limit}
    \uses{def:cvg_distribution}
Let $X_1, X_2, \ldots$ be i.i.d. real random variables with mean 0 and variance 1, and let $Z$ be a random variable with law $\mathcal N(0,1)$. Then
\begin{align*}
\frac{1}{\sqrt{n}}\sum_{k=1}^n X_k \xrightarrow{d} Z \: .
\end{align*}
\end{theorem}

\begin{proof}\uses{lem:charFun_smul, lem:charFun_add_of_indep, thm:charFun_tendsto_iff_measure_tendsto, lem:charFun_taylor, lem:gaussian_charFun, lem:tendsto_pow_exp_of_isLittleO, lem:iIndepFun_iff_pi_map_eq_map}
\leanok
Let $S_n = \frac{1}{\sqrt{n}}\sum_{k=1}^n X_k$. Let $\phi$ be the characteristic function of $X_k$. By Lemma~\ref{lem:charFun_smul} and~\ref{lem:charFun_add_of_indep}, the characteristic function of $S_n$ is $\phi_n(t) = (\phi(n^{-1/2}t))^n$.

By Lemma \ref{lem:charFun_taylor},
\begin{align*}
\phi_n(t) = (\phi(n^{-1/2}t))^n = \left(1 - \frac{1}{2n}t^2 + o(\frac{1}{n})\right)^n \to_{n \to +\infty} e^{-t^2/2} \: .
\end{align*}

Since the r.h.s. is the characteristic function of $Z$, we conclude that $\frac{1}{\sqrt{n}}\sum_{k=1}^n X_k \xrightarrow{d} Z$ by Lemma~\ref{thm:charFun_tendsto_iff_measure_tendsto}.
\end{proof}


\section{The central limit theorem for random variables in $\mathbb{R}^d$}

\begin{theorem}[Cramér-Wold]\label{thm:cramer_wold}
    \uses{def:cvg_distribution}
Let $X, X_1, X_2, \ldots$ be random variables in $\mathbb{R}^d$.
Then $X_n \xrightarrow{d} X$ iff for every $a \in \mathbb{R}^d$, $\langle a, X_n \rangle \xrightarrow{d} \langle a, X \rangle$.
\end{theorem}

\begin{proof}
    \uses{thm:charFun_tendsto_iff_measure_tendsto}
By Theorem~\ref{thm:charFun_tendsto_iff_measure_tendsto}, the convergence in distribution of the random variables is equivalent to the convergence of their characteristic functions for each $a \in \mathbb{R}^d$, which for each such $a$ is equivalent to the convergence in distribution of the random variables $\langle a, X_n \rangle$ by the same theorem.
\end{proof}

\begin{theorem}[Central limit theorem in $\mathbb{R}^d$]\label{thm:multivariate_clt}
    %\leanok
    %\lean{ProbabilityTheory.multivariate_central_limit}
Let $X_1, X_2, \ldots$ be i.i.d. random variables in $\mathbb{R}^d$ with mean 0 and identity covariance matrix, and let $Z$ be a random variable with law $\mathcal N(0,I_d)$. Then
\begin{align*}
\frac{1}{\sqrt{n}}\sum_{k=1}^n X_k \xrightarrow{d} Z \: .
\end{align*}
\end{theorem}

\begin{proof}
\uses{thm:cramer_wold, thm:clt, lem:stdGaussian_finiteDimensional}

\end{proof}

%%FILE blueprint/src/chapter/gaussian.tex
\chapter{Gaussian measures}

\begin{definition}\label{def:gaussianReal}
\mathlibok
\lean{ProbabilityTheory.gaussianReal}
A measure on $\mathbb{R}$ is Gaussian if it is equal to $\mathcal{N}(m, \sigma^2)$ for some $m \in \mathbb{R}$ and $\sigma^2 \ge 0$, in which $\mathcal{N}(m, \sigma^2)$ is the measure absolutely continuous with respect to the Lebesgue measure with density $x \mapsto \frac{1}{\sqrt{2 \pi \sigma^2}}e^{- \frac{1}{2\sigma^2}(x - m)^2}$ if $\sigma^2>0$ and the Dirac probability measure at $m$ if $\sigma^2 = 0$.
\end{definition}


\begin{lemma}\label{lem:isProbabilityMeasure_gaussianReal}
\mathlibok
\uses{def:gaussianReal}
A real Gaussian measure is a probability measure.
\end{lemma}

\begin{proof}\leanok
\end{proof}


\begin{lemma}\label{lem:gaussian_charFun}
\mathlibok
\lean{ProbabilityTheory.charFun_gaussianReal}
\uses{def:charFun, def:gaussianReal}
The Gaussian distribution $\mathcal N(m, \sigma^2)$ has characteristic function $\phi(t) = e^{itm - \sigma^2 t^2 /2}$.
\end{lemma}

\begin{proof}\leanok
\end{proof}


\begin{lemma}\label{lem:add_gaussianReal}
\mathlibok
\lean{ProbabilityTheory.gaussianReal_conv_gaussianReal}
\uses{def:gaussianReal}
The sum of two independent real Gaussian random variables is Gaussian.
\end{lemma}

\begin{proof}\leanok
\uses{lem:gaussian_charFun, lem:charFun_add_of_indep}

\end{proof}


\begin{definition}\label{def:isGaussian}
\uses{def:gaussianReal}
\mathlibok
\lean{ProbabilityTheory.IsGaussian}
A Borel measure $\mu$ on a separable Banach space $E$ is said to be Gaussian if for all continuous linear maps $\ell : E \to \mathbb{R}$, the pushforward $\ell_*\mu$ is a real Gaussian measure.
\end{definition}

\begin{lemma}\label{lem:isProbabilityMeasure_isGaussian}
\mathlibok
\uses{def:isGaussian}
A Gaussian measure is a probability measure.
\end{lemma}

\begin{proof}\leanok
\uses{lem:isProbabilityMeasure_gaussianReal}
\end{proof}


\begin{lemma}\label{lem:isGaussian_gaussianReal}
\mathlibok
\lean{ProbabilityTheory.isGaussian_gaussianReal}
\uses{def:gaussianReal, def:isGaussian}
The real Gaussian measures are Gaussian measures in the sense of Definition~\ref{def:isGaussian}.
\end{lemma}

\begin{proof}\leanok

\end{proof}


\begin{lemma}\label{lem:isGaussian_add}
\mathlibok
\lean{ProbabilityTheory.isGaussian_conv}
\uses{def:isGaussian}
The sum of two independent Gaussian random variables is Gaussian.
\end{lemma}

\begin{proof}\leanok
\uses{lem:add_gaussianReal,lem:isProbabilityMeasure_isGaussian}

\end{proof}


\begin{lemma}\label{lem:stdGaussian_finiteDimensional}
\uses{def:isGaussian}
Let $E$ be a finite dimensional real inner product space and let $b_1, \ldots, b_d$ be an orthonormal basis of $E$.
Let $X_1, \ldots, X_d$ be independent standard Gaussian random variables on $\mathbb{R}$.
Then the law of $X_1 b_1 + \ldots + X_d b_d$ is a Gaussian measure on $E$.
\end{lemma}

\begin{proof}
\uses{lem:isGaussian_add}

\end{proof}

%%FILE blueprint/src/chapter/main.tex
% This is the main point of entry to the blueprint.
% Add chapters of the blueprint here.
% This file is not meant to be built. Build src/web.tex or src/print.text instead.

This document is a blueprint for the formalization of the central limit theorem in Lean.
For the proof, we will introduce and use characteristic functions, as well as the notion of tightness of sets of measures.

The main references for the results detailed here are
\begin{itemize}
	\item chapter 6 of Foundations of Modern Probability (3rd ed.), by Olav Kallenberg,
	\item chapter 8 of Measure Theory, volume II, by V. Bogachev.
	\item chapter 3 of Markov Processes, Characterization and Convergence, by S.N. Ethier and T.G. Kurtz.
\end{itemize}

This blueprint is maintained by R\'emy Degenne. Part of the proof is inspired by a document written by Peter Pfaffelhuber, in particular the use in the proof of the separating sigma-algebra of exponential polynomials as in the book by Ethier and Kurtz.

\input{chapter/mathlib_defs}
\input{chapter/tight}
\input{chapter/separating}
\input{chapter/char_fun}
\input{chapter/gaussian}
\input{chapter/clt}


%%FILE blueprint/src/chapter/mathlib_defs.tex
\chapter{Convergence definitions}

% \begin{definition}[Convergence in probability]\label{def:cvg_probability}
% \mathlibok
% \lean{MeasureTheory.TendstoInMeasure}
% We say that $(X_i)_{i \in \mathbb{N}}$ tends to $X$ in probability (or in measure) in a space with distance function $d$ if for all $\varepsilon > 0$
% \begin{align*}
% \mathbb{P}(d(X_n, X) \ge \varepsilon) \to_{n \to +\infty} 0
% \end{align*}
% We write $X_n \xrightarrow{p} X$.
% \end{definition}

\begin{definition}[Weak convergence of probability measures]\label{def:weak_cvg_measure}
\mathlibok
\lean{MeasureTheory.ProbabilityMeasure.tendsto_iff_forall_integral_tendsto}
In a topological space $S$ in which the opens are measurable, we say that a sequence of probability measures $\mu_n$ converges weakly to a probability measure $\mu$ and write $\mu_n \xrightarrow{w} \mu$ if $\mu_n[f] \to \mu[f]$ for every bounded continuous function $f : S \to \mathbb{R}$.
\end{definition}

We write $\mathcal L(X)$ for the law of a random variable $X : \Omega \to S$. This is a measure on $S$, defined as the map of the measure on $\Omega$ by $X$.

\begin{definition}[Convergence in distribution]\label{def:cvg_distribution}
\uses{def:weak_cvg_measure}
\mathlibok
\lean{MeasureTheory.ProbabilityMeasure.tendsto_iff_forall_integral_tendsto}
We say that $X_n$ converges in distribution to $X$ and write $X_n \xrightarrow{d} X$ if $\mathcal L(X_n) \xrightarrow{w} \mathcal L(X)$.
\end{definition}

%%FILE blueprint/src/chapter/separating.tex
\chapter{Separating subalgebras}

\begin{definition}\label{def:separates_points}
\mathlibok
\lean{Set.SeparatesPoints}
A set $\mathcal F$ of functions $E \to F$ separates points in $E$ if for all $x, y \in E$ with $x \ne y$, there exists $f \in \mathcal F$ with $f(x) \ne f(y)$.
\end{definition}


\begin{lemma}\label{lem:bounded_continuous_separating}
\mathlibok
\lean{MeasureTheory.ext_of_forall_integral_eq_of_IsFiniteMeasure}
In a Borel space $E$, if two probability measures $\mu, \nu$ on $E$ are such that $\mu[f] = \nu[f]$ for all $f \in C_b(E, \mathbb{C})$, the set of bounded continuous functions from $E$ to $\mathbb{C}$, then $\mu = \nu$.
\end{lemma}

\begin{proof}\leanok
The Mathlib lemma \texttt{MeasureTheory.FiniteMeasure.ext\_of\_forall\_lintegral\_eq} shows that the result is true for $C_b(E, \mathbb{R}_+)$. Since $C_b(E, \mathbb{R}_{+}) \subseteq C_b(E, \mathbb{C})$, we get the result.
\end{proof}


\begin{lemma}\label{lem:exp_character}
\mathlibok
\lean{BoundedContinuousFunction.charMonoidHom}
$x \mapsto (t \mapsto \exp(i \langle t, x \rangle))$ is a monoid homomorphism from $(E,+)$ to $C_b(E, \mathbb{C})$.
\end{lemma}

\begin{proof}\leanok
\end{proof}


\begin{lemma}\label{lem:starSubalgebra_expPoly}
\mathlibok
\lean{BoundedContinuousFunction.charPoly}
The functions of the form $x \mapsto \sum_{k=1}^n a_k e^{i\langle t_k, x\rangle}$ for $n \in \mathbb{N}$, $a_1, \ldots, a_n \in \mathbb{R}$ and $t_1, \ldots, t_n \in E$ are a star-subalgebra of $C_b(E, \mathbb{C})$. 
\end{lemma}

\begin{proof}\uses{lem:exp_character} \leanok
The monoid homomorphism can be lifted to a homomorphism of algebras from \texttt{AddMonoidAlgebra ℝ E} to $C_b(E, \mathbb{C})$ using \texttt{AddMonoidAlgebra.lift}. $\mathcal M$ is the range of this homomorphism. The ``star'' part can be checked easily.
\end{proof}


Let $\mathcal M$ be the set of these functions, which we call exponential polynomials.

\begin{lemma}\label{lem:separates_points_expPoly}
\uses{def:separates_points, lem:starSubalgebra_expPoly}
\mathlibok
\lean{BoundedContinuousFunction.separatesPoints_charPoly}
The star-subalgebra $\mathcal M$ separates points.
\end{lemma}

\begin{proof}\leanok
\end{proof}


\begin{lemma}\label{lem:integral_restrict_compact}
\mathlibok
\lean{abs_integral_sub_setIntegral_mulExpNegMulSq_comp_lt}
Let $\mu$ be a finite measure on a Borel space $E$ and let $K$ be a measurable compact set such that $\mu(K^c) \le \varepsilon$. Let $f \in C(E, \mathbb{R})$. Then
\begin{align*}
\left\vert \mu[fe^{-\varepsilon \Vert f \Vert^2}] - \mu[f e^{-\varepsilon \Vert f \Vert^2} \mathbb{I}_K] \right\vert
\le \sqrt{\varepsilon} \: .
\end{align*}
\end{lemma}

\begin{proof}\leanok
\end{proof}


\begin{lemma}\label{lem:introduce_exponential}
\mathlibok
\lean{tendsto_integral_mulExpNegMulSq_comp}
For any $f \in C_b(E, \mathbb{R})$ and a probability measure $\mu$,
\begin{align*}
\mu[f] 
= \lim_{\varepsilon \to 0} \mu\left[f e^{-\varepsilon \Vert f \Vert^2} \right]
\: .
\end{align*}
\end{lemma}

\begin{proof}\leanok
\end{proof}


\begin{lemma}\label{lem:dist_integral_mulExpNegMulSq_comp_le}
\uses{def:separates_points}
\mathlibok
\lean{dist_integral_mulExpNegMulSq_comp_le}
Let $\mu, \nu$ be two finite measures in a standard Borel space $E$. Let $\mathcal{A}$ be a subalgebra of $C_b(E, \mathbb{R})$ that separates points.
If for all $g \in \mathcal A$, $\mu[g] = \nu[g]$, then for all $f \in C_b(E, \mathbb{R})$ and all $\varepsilon > 0$,
\begin{align*}
    \left\vert \mu[fe^{-\varepsilon \Vert f \Vert^2}] - \nu[f e^{-\varepsilon \Vert f \Vert^2}] \right\vert
    \le 6\sqrt{\varepsilon} \: .
\end{align*}
\end{lemma}

\begin{proof}\leanok
\uses{lem:integral_restrict_compact}
\end{proof}


\begin{theorem}[Subalgebras separating points]\label{thm:separating_starSubalgebra}
\mathlibok
\lean{MeasureTheory.ext_of_forall_mem_subalgebra_integral_eq_of_polish}
Let $E$ be a complete separable pseudo-metric space. Let $\mathcal A \subseteq C_b(E, \mathbb{C})$ be a star-subalgebra that separates points.
Then if two probability measures $\mu, \nu$ of $E$ are such that $\mu[f] = \nu[f]$ for all $f \in \mathcal A$, then $\mu = \nu$.
\end{theorem}

\begin{proof}\uses{lem:bounded_continuous_separating, lem:introduce_exponential, lem:dist_integral_mulExpNegMulSq_comp_le}
\leanok
Let $\mu$ and $\nu$ be two probability measures such that $\mu[f] = \nu[f]$ for all $f \in \mathcal A$. We want to prove that $\mu = \nu$.
By Lemma~\ref{lem:bounded_continuous_separating}, it suffices to prove that for all $g \in C_b(E, \mathbb{C})$, $\mu[g] = \nu[g]$.

This is proved in Mathlib PR \#19782.
\end{proof}


\begin{lemma}\label{lem:cvg_of_separating}
\leanok
\lean{MeasureTheory.ProbabilityMeasure.tendsto_of_tight_of_separatesPoints}
\uses{def:weak_cvg_measure}
Let $\mathcal A$ be a subalgebra of $C_b(E, \mathbb{C})$ that separates points. Let $\mu, \mu_1, \mu_2, \ldots$ be probability measures, with $(\mu_n)$ tight. If for all $f \in \mathcal A$, $\mu_n[f] \to \mu[f]$, then $\mu_n \xrightarrow{w} \mu$.
\end{lemma}

\begin{proof}\uses{lem:relatively_compact_of_tight, thm:separating_starSubalgebra}
\leanok
To show convergence to $\mu$, it suffices to show that every subsequence has a subsequence that converges to $\mu$ (\texttt{Filter.tendsto\_of\_subseq\_tendsto}).
Since the family is tight, it is relatively compact by Prokhorov's theorem. We get that each subsequence has a converging subsequence.
Let then $\mu'$ be a weak limit of $\mu_n$. We have $\mu_n[f] \to \mu'[f]$ and $\mu_n[f] \to \mu[f]$ for all $f \in \mathcal A$, hence $\mu'[f] = \mu[f]$ for all $f \in \mathcal A$. Since $\mathcal A$ separates points, $\mathcal \mu' = \mu$ by Theorem~\ref{thm:separating_starSubalgebra}.
\end{proof}


%%FILE blueprint/src/chapter/tight.tex
\chapter{Tight families of measures}

\begin{definition}\label{def:tight}
\lean{MeasureTheory.IsTightMeasureSet} \mathlibok
A set $S$ of measures on $\Omega$ is tight if for all $\varepsilon > 0$ there exists a compact set $K$ such that for all $\mu \in S$, $\mu(K^c) \le \varepsilon$.
\end{definition}


\begin{lemma}\label{lem:isTightMeasureSet_iff_norm}
\uses{def:tight}
\mathlibok
\lean{MeasureTheory.isTightMeasureSet_iff_tendsto_measure_norm_gt}
Let $S$ be a set of measures on a finite-dimensional normed space $E$. 
Then $S$ is tight iff $\lim_{r \to +\infty}\sup_{\mu \in S} \mu\{\Vert x \Vert > r\} = 0$.
\end{lemma}

\begin{proof}\leanok

\end{proof}
	

\begin{lemma}\label{lem:isTightMeasureSet_iff_sup_inner}
\uses{def:tight}
\mathlibok
\lean{MeasureTheory.isTightMeasureSet_iff_inner_tendsto}
Let $S$ be a set of measures on a finite-dimensional inner product space $E$.
Then $S$ is tight iff for all $y \in E$, $\lim_{r \to +\infty}\sup_{\mu \in S} \mu\{\vert\langle y, x\rangle\vert > r\} = 0$.
\end{lemma}

\begin{proof}\leanok
\uses{lem:isTightMeasureSet_iff_norm}
\end{proof}


\begin{lemma}\label{lem:isTightMeasureSet_iff_basis}
\uses{def:tight}
\leanok
\lean{isTightMeasureSet_iff_tendsto_limsup_inner}
Let $(\mu_n)_{n \in \mathbb{N}}$ be a sequence of measures on a finite-dimensional inner product space $E$. 
Then the set $\{\mu_n \mid n \in \mathbb{N}\}$ is tight iff for all $y \in E$, $\lim_{r \to +\infty}\limsup_{n \to \infty} \mu_n\{\vert\langle y, x\rangle\vert > r\} = 0$.
\end{lemma}

\begin{proof}\leanok
\uses{lem:isTightMeasureSet_iff_sup_inner}
\end{proof}


% \begin{lemma}\label{lem:tight_of_cvg}
% \uses{def:tight}
% If $\mu_1, \mu_2, \ldots$ converge weakly to $\mu$, then $\{\mu_n \mid n \in \mathbb{N}\}$ is tight.
% \end{lemma}

% \begin{proof}
% Fix $r > 0$ and let $f(x) = (1 - (r - \vert x \vert)_+)_+$ . $f$ is a bounded continuous function. Then
% \begin{align*}
% \limsup_{n \to +\infty} \mu_n\{\vert X \vert > r\}
% \le \lim_{n \to +\infty} \mu_n\left[ f \right]
% = \mu[f]
% \le \mu\{\vert X \vert > r - 1\} \: .
% \end{align*}
% As $r \to +\infty$, the last quantity tends to 0.
% \end{proof}


% \begin{lemma}\label{tight_iff_cvg_probability}
% \uses{def:tight, def:cvg_probability}
% For random variables $(X_i)_{i \in \mathbb{N}}$ (TODO in which kind of space?), the two following conditions are equivalent:
% \begin{itemize}
% 	\item $(\mathcal L(X_n))$ is tight,
% 	\item For all $(c_n) \ge 0$ with $c_n \to 0$, $c_n X_n \xrightarrow{p} 0$.
% \end{itemize}
% \end{lemma}

% \begin{proof}

% \end{proof}


% \begin{lemma}\label{tight_closure}
% \uses{def:tight}
% The closure of a tight set of measures is tight.
% \end{lemma}

% \begin{proof}
% \end{proof}



\section{Prokhorov's theorem}

% \begin{lemma}\label{lem:exists_finite_union_inter_lt}
% Let $(U_n)_{n \in \mathbb{N}}$ be open sets in a complete separable metric space $E$ such that $\bigcup_{n=1}^{+ \infty} U_n = E$. Let $\Gamma$ be a relatively compact set of probability measures on $E$ for the topology of weak convergence of measures. Then for all $\varepsilon > 0$ there exists a finite set $S \subseteq \mathbb{N}$ such that for all $\mu \in \Gamma$, $\mu(\bigcup_{n \in S} U_n) > 1 - \varepsilon$.
% \end{lemma}

% \begin{proof}
% \end{proof}

% \begin{lemma}\label{lem:prokhorov_aux1}
% \uses{def:tight}
% Let $E$ be a complete separable metric space and let $S \subseteq \mathcal P(E)$. If the closure of $S$ is compact, then $S$ is tight.
% \end{lemma}

% \begin{proof}\uses{lem:exists_finite_union_inter_lt}
% \end{proof}

% \begin{theorem}[Riesz-Markov-Kakutani representation theorem, for compact spaces]\label{thm:riesz_markov_kakutani}
% Let $E$ be a compact T2 space and let $\psi$ be a positive linear functional on $C(E, \mathbb{C})$. There exists a regular measure $\mu$ finite on compacts such that
% \begin{align*}
% \forall f \in C(E, \mathbb{C}), \ \psi(f) = \int_{x \in E} f(x) \partial\mu(x) \: .
% \end{align*}
% \end{theorem}

% \begin{proof}
% \end{proof}

% \begin{lemma}\label{lem:ProbabilityMeasure_compact}
% If $E$ is a compact metric space, then $\mathcal P(E)$ with the Prokhorov metric is a compact metric space.
% \end{lemma}

% \begin{proof}\uses{thm:riesz_markov_kakutani}
% \end{proof}

\begin{theorem}[Prokhorov's theorem]\label{thm:prokhorov}
\notready
\lean{isTightMeasureSet_iff_isCompact_closure}
Let $E$ be a complete separable metric space and let $S \subseteq \mathcal P(E)$. Then the following are equivalent:
\begin{enumerate}
	\item $S$ is tight.
	\item the closure of $S$ is compact.
\end{enumerate}
\end{theorem}

\begin{proof}
\end{proof}


\begin{lemma}\label{lem:relatively_compact_of_tight}
\leanok
\lean{isCompact_closure_of_isTightMeasureSet}
Let $E$ be a metric space and let $S \subseteq \mathcal P(E)$ be tight. Then the closure of $S$ is compact.
\end{lemma}

\begin{proof}\uses{thm:prokhorov}
\end{proof}

